\section{Metric Definitions}
\label{app:metrics}

\begin{table}[h]
\centering
\caption{Detailed metric dictionary for benchmark logs.}
\label{tab:metric-dictionary}
\small
\begin{tabularx}{\linewidth}{@{}lX@{}}
\toprule
\textbf{Metric} & \textbf{Definition} \\
\midrule
\texttt{input\_tokens} & Total input tokens reported by the provider for one model call. \\
\texttt{cached\_input\_tokens} & Input tokens reported as cached by the provider. If absent, mark as missing rather than assuming zero. \\
\texttt{uncached\_input\_tokens} & \texttt{input\_tokens - cached\_input\_tokens}. This is the primary paid-input reduction target. \\
\texttt{cache\_hit\_rate} & \texttt{cached\_input\_tokens / input\_tokens}. \\
\texttt{estimated\_actual\_cost} & Cost using uncached input price, cached input price, and output price. \\
\texttt{estimated\_full\_cost} & Counterfactual cost if all input tokens were uncached. \\
\texttt{ttft\_ms} & Time to first token, measured at the client or trace layer. \\
\texttt{task\_success} & Boolean or reviewer-scored task outcome using a fixed rubric. \\
\texttt{validation\_passed} & Whether the task-specific validation command passed. \\
\bottomrule
\end{tabularx}
\end{table}

\section{Interpretation Rules}
\label{app:interpretation}

A run should not be considered successful merely because it reports lower total tokens. The preferred claim is narrower and more honest:
\begin{quote}
The cache-friendly condition reduced paid uncached input while preserving task success.
\end{quote}
If a strategy lowers cost but increases failed tasks, extra turns, or unintended file changes, it should be reported as a quality regression rather than a win.

